\section{misc}

\subsection{pv}

IEEEeqnarry, multi no., no good 
\begin{IEEEeqnarray}{r}
	pv(i, j) = \left\{ 
	\begin{array}{rl}
		max(x_{i - 1, j}, x_{i, j - 1})                & , x_{i - 1, j - 1} \leq min(x_{i - 1, j}, x_{i, j - 1}) \\
		min(x_{i - 1, j}, x_{i, j - 1})                & , x_{i - 1, j - 1} \geq max(x_{i - 1, j}, x_{i, j - 1}) \\
		x_{i - 1, j} + x_{i, j - 1} - x_{i - 1, j - 1} & ,  \textrm{otherwise}
	\end{array} \right. \\
	, \textrm{where } 1 \leq i < \textrm{image height}, 1 \leq j < \textrm{image width} 
\end{IEEEeqnarray}

\rule{\linewidth}{1px}

equation with array of array, might be okay. 
\begin{equation}
	\begin{array}{l}
%		\forall i \in (1, h], \forall j \in (1, w], \\
		pv(i, j) = \left\{ 
		\begin{array}{rl}
			max(x_{i - 1, j}, x_{i, j - 1})                & \text{if } x_{i - 1, j - 1} \leq min(x_{i - 1, j}, x_{i, j - 1}), \\
			min(x_{i - 1, j}, x_{i, j - 1})                & \text{if } x_{i - 1, j - 1} \geq max(x_{i - 1, j}, x_{i, j - 1}), \\
			x_{i - 1, j} + x_{i, j - 1} - x_{i - 1, j - 1} & \text{otherwise}
		\end{array} \right.
		\\
		, \textrm{where } 1 \leq i < \textrm{image height}, 1 \leq j < \textrm{image width}
	\end{array}
\end{equation}

equation with array
\begin{equation}
	pv(i, j) = \left\{ 
	\begin{array}{rl}
		max(x_{i - 1, j}, x_{i, j - 1})                & , x_{i - 1, j - 1} \leq min(x_{i - 1, j}, x_{i, j - 1}) \\
		min(x_{i - 1, j}, x_{i, j - 1})                & , x_{i - 1, j - 1} \geq max(x_{i - 1, j}, x_{i, j - 1}) \\
		x_{i - 1, j} + x_{i, j - 1} - x_{i - 1, j - 1} & ,  \textrm{otherwise}
	\end{array} \right.
\end{equation}

equation, array of array, seems good.
\begin{equation}
	\begin{array}{r}
		pv(i, j) = \left\{ 
		\begin{array}{rl}
			max(x_{i - 1, j}, x_{i, j - 1})                & , x_{i - 1, j - 1} \leq min(x_{i - 1, j}, x_{i, j - 1}) \\
			min(x_{i - 1, j}, x_{i, j - 1})                & , x_{i - 1, j - 1} \geq max(x_{i - 1, j}, x_{i, j - 1}) \\
			x_{i - 1, j} + x_{i, j - 1} - x_{i - 1, j - 1} & ,  \textrm{otherwise}
		\end{array} \right.
		\\
		, \textrm{where } 1 \leq i < \textrm{image height}, 1 \leq j < \textrm{image width}
	\end{array}
\end{equation}

equation with cases, also okay
\begin{equation}
	pv(i, j) = 
	\begin{cases}
		max(x_{i - 1, j}, x_{i, j - 1})                ,& x_{i - 1, j - 1} \leq min(x_{i - 1, j}, x_{i, j - 1}) \\
		min(x_{i - 1, j}, x_{i, j - 1})                ,& x_{i - 1, j - 1} \geq max(x_{i - 1, j}, x_{i, j - 1}) \\
		x_{i - 1, j} + x_{i, j - 1} - x_{i - 1, j - 1} ,& \textrm{otherwise}
	\end{cases} 
\end{equation}

\rule{\linewidth}{1px}

multline with array, this seems fine.
\begin{multline}
	\forall i \in [1, h), \forall j \in [1, w), \\
	pv(i, j) = \left\{ 
	\begin{array}{rl}
		max(x_{i - 1, j}, x_{i, j - 1})                & \text{if } x_{i - 1, j - 1} \leq min(x_{i - 1, j}, x_{i, j - 1}), \\
		min(x_{i - 1, j}, x_{i, j - 1})                & \text{if } x_{i - 1, j - 1} \geq max(x_{i - 1, j}, x_{i, j - 1}), \\
		x_{i - 1, j} + x_{i, j - 1} - x_{i - 1, j - 1} & \text{otherwise}
	\end{array} \right.
% \\ , \text{where } h \text{ and } w \text{ denote the height and the weight of the image, respectively.}
\end{multline}
, where  $h$  and  $w$  denote the height and the weight of the image, respectively.

multline with cases, best so far.

\begin{multline}
	pv(i, j) =  
	\begin{cases}
		max(x_{i - 1, j}, x_{i, j - 1})               , &  \text{if } x_{i - 1, j - 1} \leq min(x_{i - 1, j}, x_{i, j - 1}), \\
		min(x_{i - 1, j}, x_{i, j - 1})               , &  \text{if } x_{i - 1, j - 1} \geq max(x_{i - 1, j}, x_{i, j - 1}), \\
		x_{i - 1, j} + x_{i, j - 1} - x_{i - 1, j - 1}, &  \textrm{otherwise}. 
	\end{cases} 
%	\textrm{where } 1 \leq i < n, 1 \leq j < m
\end{multline}



\rule{\linewidth}{1px}

multline, this is just too much
\begin{multline}
	pv(i, j) =  
	\begin{cases}
		max(x_{i - 1, j}, x_{i, j - 1})               , &  \text{if } x_{i - 1, j - 1} \leq min(x_{i - 1, j}, x_{i, j - 1}), \\
		min(x_{i - 1, j}, x_{i, j - 1})               , &  \text{if } x_{i - 1, j - 1} \geq max(x_{i - 1, j}, x_{i, j - 1}), \\
		x_{i - 1, j} + x_{i, j - 1} - x_{i - 1, j - 1}, &  \textrm{otherwise}, 
	\end{cases} \\
	\textrm{and } \\
	i = \{x \mid x \in \mathbb{N}, 1 \leq x < n\} \\
	j = \{y \mid y \in \mathbb{N}, 1 \leq y < m\} \\
	\textrm{where } n \text{ is the height of image and} \\
	m \text{ is the width of image} 
\end{multline}

multline with case, seems okay
\begin{multline}
	pv(i, j) =  
	\begin{cases}
		max(x_{i - 1, j}, x_{i, j - 1})               , &  \text{if } x_{i - 1, j - 1} \leq min(x_{i - 1, j}, x_{i, j - 1}), \\
		min(x_{i - 1, j}, x_{i, j - 1})               , &  \text{if } x_{i - 1, j - 1} \geq max(x_{i - 1, j}, x_{i, j - 1}), \\
		x_{i - 1, j} + x_{i, j - 1} - x_{i - 1, j - 1}, &  \textrm{otherwise}, 
	\end{cases} 
\end{multline}
where
\begin{eqnarray}
	i = \{x \mid x \in \mathbb{N}, 1 \leq x < n\} \notag \\
	j = \{y \mid y \in \mathbb{N}, 1 \leq y < m\} \notag \\
	n \text{ is the height of image} \notag \\
	m \text{ is the width of image} \notag
\end{eqnarray}

align with cases and intertext, best so far, but still too much. \\
the where clause are aligned but the notag thing is kind of annoying.
\begin{align}
	pv(i, j) &=  
		\begin{cases}
			max(x_{i - 1, j}, x_{i, j - 1})               , &  \text{if } x_{i - 1, j - 1} \leq min(x_{i - 1, j}, x_{i, j - 1}), \\
			min(x_{i - 1, j}, x_{i, j - 1})               , &  \text{if } x_{i - 1, j - 1} \geq max(x_{i - 1, j}, x_{i, j - 1}), \\
			x_{i - 1, j} + x_{i, j - 1} - x_{i - 1, j - 1}, &  \textrm{otherwise}, 
		\end{cases}
	\intertext{where}
	i &= \{x \mid x \in \mathbb{N}, 1 \leq x < n\} \notag \\
	j &= \{y \mid y \in \mathbb{N}, 1 \leq y < m\} \notag \\
	n &\text{ is the height of image} \notag \\
	m &\text{ is the width of image} \notag
\end{align}


\rule{\linewidth}{1px}









\subsection{pe}

\begin{equation}
	pe(i, j) = pv(i, j) \oplus pixel(i, j), \textrm{where } 1 \leq i < \textrm{image height}, 1 \leq j < \textrm{image width}
\end{equation}

\begin{eqnarray}
	pe(i, j) = pv(i, j) \oplus pixel(i, j), \\
	\textrm{where } 1 \leq i < \textrm{image height}, 1 \leq j < \textrm{image width}
\end{eqnarray}


\begin{multline}
	pe(i, j) = pv(i, j) \oplus pixel(i, j), \\
	\textrm{where } 1 \leq i < \textrm{image height}, 1 \leq j < \textrm{image width}
\end{multline}

\begin{equation}
	pe(i, j) = pv(i, j) \oplus pixel(i, j), \forall 1 \leq i < n, 1 \leq j < m
\end{equation}

\begin{align}
	pe(i, j) & = & pv(i, j) \oplus pixel(i, j),  \\
	& & \textrm{where } 1 \leq i < \textrm{image height}, 1 \leq j < \textrm{image width}
\end{align}





\subsection{label map t}

\begin{equation}
	label(i, j) = \left\{ 
	\begin{array}{ll}
		8                       & , pe(i, j) = 0       \\
		7 - log_2(pe(i, j)) + 1 & , \textrm{otherwise}
	\end{array} \right.
\end{equation}

\begin{equation}
	label(i, j) = 
	\begin{cases}
		8                       & , pe(i, j) = 0       \\
		7 - log_2(pe(i, j)) + 1 & , \textrm{otherwise}
	\end{cases} 
\end{equation}






\subsection{length of label map}

\begin{equation}
	LM = \sum_{t=0}^{8} (freq[t] \times \textrm{the length of Huffman code of }t)
\end{equation}

\begin{equation}
	LM = \sum_{t=0}^{8} (freq[t] \cdot \textrm{the length of Huffman code of }t)
\end{equation}

\begin{multline}
	LM = \sum_{t=0}^{8} (freq(t) \cdot \lvert HuffEncode(t) \rvert) , \text{ where } \\
	HuffEncode(x) \text{ denotes the Huffman code of } x \text{ and } \\ 
	\lvert y \rvert \text{ denotes the length of } y
\end{multline}

\begin{equation}
	\begin{aligned}
		LM & = \sum_{t=0}^{8} (freq(t) \cdot \lvert HuffEncode(t) \rvert) , \\
		& \text{ where } \\
		& HuffEncode(x) \text{ denotes the Huffman code of } x \text{ and } \\ 
		& \lvert y \rvert \text{ denotes the length of } y
	\end{aligned}
\end{equation}





\subsection{embed}
\begin{equation}
	{x_e}\prime(i, j) = \left\{
	\begin{array}{rl}
		x_e(i, j) \bmod {2^{7-t}} + \sum_{s=0}^t(b_s \times 2^{7-s}) & , 0 \leq t < 7    \\
		\sum_{s=1}^8(b_s \times 2^{8-s})                            & , 7 \leq t \leq 8
	\end{array}
	\right.
\end{equation}



\subsection{extract}

\begin{equation}
	x(i, j) = \left\{
	\begin{array}{ll}
		px(i, j) / {2^{7-t}} + b_{t+1} \times 2^{7-t} + x_{ew}(i, j) \bmod {2^{7-t}} & , 0 \leq t < 8    \\
		px(i, j)                            & ,  t = 8
	\end{array}
	\right.
\end{equation}




\subsection{misc}

\begin{equation}
	\begin{array}{rcl}
		MSB(x, b) & = & x ~\&~  (((1 \ll b) - 1) \ll (8 - b))  \\
		LSB(x, b) & = & x ~\&~  ((1 \ll b) - 1)
	\end{array}
\end{equation}

\begin{table}
	\caption{text}
	\centering
	\begin{tabular}{|c | c |}
		\hline 
		$x_1$ & $x_2$ \\
		\hline 
		$x_3$ & $x$   \\
		\hline 
	\end{tabular}
\end{table}








\begin{figure}
	\centering
	\begin{tabular}{|p{.25 \textwidth} | p{.25 \textwidth} |p{.25 \textwidth} | p{.25 \textwidth} |}
		\hline 
		32 bits of Huffman coding rule & 20 bits of the length of the Huffman encoded label map & Huffman encoded label map & 1023 bytes of reference bytes \\ 
		\hline
	\end{tabular}
	\caption{text}
\end{figure}

\begin{figure}
	\centering
	\begin{tabular}{|p{.20 \textwidth} | p{.20 \textwidth} |p{.20 \textwidth} | p{.20 \textwidth}| p{.20 \textwidth} |}
		\hline 
		17 bits of Huffman coding tree topology & 36 bits of Huffman coding value & 20 bits of the length of the Huffman encoded label map & Huffman encoded label map & 16 bytes of reference bytes \\ 
		\hline
	\end{tabular}
	\caption{text}
\end{figure}





\begin{figure}
	\centering
	
	\begin{tabular}{rcl}
		$x_e$ & $=$     & 
		\begin{tabular}{| c | c | c | c | c | c | c | c |}
			\hline
			\cellcolor{lightgray}0 & \cellcolor{lightgray}0 & \cellcolor{lightgray}1 & \cellcolor{gray}1 & 0 & 1 & 0 & 1 \\
			\hline
		\end{tabular} 
		\\
		& &
		\begin{tabular}{ c  c  c  c  c  c  c  c }
			& $t = 3$ & &  &  & &  & \\
		\end{tabular} 
		\\
		secret bits    & $=$ &
		\begin{tabular}{| c | c | c | c | c  c  c  c }
			\cline{1-4}
			1 & 1 & 0 & 1 & & & & \\
			\cline{1-4}
		\end{tabular} 
		\\
		\phantom{0}      &  &       
		\\
		${x_e}\prime $ &  $=$ &
		\begin{tabular}{| c | c | c | c | c | c | c | c |}
			\hline
			\cellcolor{lightgray}1 & \cellcolor{lightgray}1 & \cellcolor{lightgray}0 & \cellcolor{gray}1 & 0 & 1 & 0 & 1 \\
			\hline
		\end{tabular}
	\end{tabular}
	
	\caption{text}
\end{figure}

\begin{figure}
	\centering
	\begin{multline*}
		x_e  = \underbrace{001}_{t = 3}\underline{1}{0101} \\
		x_e = \underbracket{001}_{t=3}\underline{1}{0101}
	\end{multline*}
	
	
	\begin{align}
		x_e & = \underbrace{001}_{t = 3}\underline{1}{0101} \\
		x_e &= \underbracket{001}_{t=3}\underline{1}{0101}
	\end{align}
	
	
	
	
	
	\caption{text}
\end{figure}